\chapter*{Вступ}
\addcontentsline{toc}{chapter}{Вступ}

\textbf{Актуальність роботи.}
На даний час карта безхмарного покриття Землі є важливим
елементом в регіональній та глобальній системах, які відстежують зміни територій,
спричинені природними та техногенними факторами.

Наприклад, оцінка шкоди, заподіяної лісовими пожежами, спричиненими вирубкою
лісів, виверження вулканів, повеней або дослідження в агрокультурі, це насамперед
моніторинг територій, сівозміни та інше. Для цього саме й потрібно мати безхмарний
композит, щоб провести дослідження. Також такі дані будуть корисними для багатьох
картографічних веб сервісів, для яких важливо мати саме останні знімки потрібних
областей. Безхмарними знімками можуть цікавитись уряди країн, наприклад для
відслідковування лісового покриву, чи детектингу інших об’єктів.

Є також інші області де потрібні такі дані, це перевезення, виробництво,
лісництво, наука.

Оскільки якість знімків Землі з кожним поколінням супутників стає все більш кращою,
то в майбутньому це дасть ще кращий результат розробленого алгоритму з вилучення
хмар.

\textbf{Мета і завдання дослідження.}

\textit{Об'єкт дослідження}~---~супутникові дані Sentinel-2 з веб-сервісу scihub
copernicus.

\textit{Предмет дослідження}~---~морфологічний оператор розширення, регресійні методи
машинного навчання.

Головною метою роботи стала побудова моделі, за допомогою якої буде
здійснюватись якісне вилучення хмарних пікселів із супутникових знімків з певної
області України.

\textit{Завдання дослідження:}
\begin{enumerate}
	\item ознайомитись з можливостями Open Data Cube;

	\item вивчити та проаналізувати вхідні дані, а саме знімки та маски хмарності;

	\item зробити аналіз алгоритмів машинного навчання для відновлення супутникових
		даних;

	\item побудувати модель вилучення хмар;

	\item прибрати візуальну різницю в кольорах між сусідніми областями супутникових
		знімків;

	\item зробити оцінку точності.
\end{enumerate}

\textbf{Наукова новизна одержаних результатів.}

Було модернізовано маску хмарності, побудована власна модель вилучення хмарних пікселів
та зроблене відновлення супутникових даних.

\textbf{Практичне значення одержаних результатів.}

Практична значущість цієї роботи є досить високою, оскільки існує потреба
організацій та приватних осіб у якісних безхмарних зображеннях територій.

Цільовою аудиторією, для якої буде корисною ця технологія, є:

\begin{enumerate}
	\item \textbf{вебсервіси}, що працюють із географічними та супутниковими даними
		(наприклад, сервіси моніторингу рослинного покриву);

	\item \textbf{агропромисловий сектор}, якому необхідні дані для дослідження стану
		сільськогосподарських угідь;

	\item \textbf{підприємства лісового господарства}, які можуть використовувати дані
		для визначення зон насадження нових дерев та загального моніторингу лісових масивів;

	\item \textbf{державні органи}, що мають потребу в моніторингу й аналізі міських
		територій (наприклад, для відстеження урбанізації та розбудови районів підрозділами
		інфраструктури). Крім того, Міністерству внутрішніх справ такі дані будуть
		корисні для виявлення осередків незаконного видобутку бурштину;

	\item \textbf{міжнародні організації}, які займаються глобальними екологічними
		проблемами, зокрема скупченням сміття в океанах або таненням льодовиків. Безхмарні
		знімки допоможуть аналізувати динаміку цих процесів;

	\item \textbf{водні та комунальні служби}, які зможуть забезпечувати збалансоване
		управління водними ресурсами конкретного регіону.
\end{enumerate}
